% Options for packages loaded elsewhere
% Options for packages loaded elsewhere
\PassOptionsToPackage{unicode}{hyperref}
\PassOptionsToPackage{hyphens}{url}
\PassOptionsToPackage{dvipsnames,svgnames,x11names}{xcolor}
%
\documentclass[
  11pt]{article}
\usepackage{xcolor}
\usepackage{amsmath,amssymb}
\setcounter{secnumdepth}{5}
\usepackage{iftex}
\ifPDFTeX
  \usepackage[T1]{fontenc}
  \usepackage[utf8]{inputenc}
  \usepackage{textcomp} % provide euro and other symbols
\else % if luatex or xetex
  \usepackage{unicode-math} % this also loads fontspec
  \defaultfontfeatures{Scale=MatchLowercase}
  \defaultfontfeatures[\rmfamily]{Ligatures=TeX,Scale=1}
\fi
\usepackage{lmodern}
\ifPDFTeX\else
  % xetex/luatex font selection
\fi
% Use upquote if available, for straight quotes in verbatim environments
\IfFileExists{upquote.sty}{\usepackage{upquote}}{}
\IfFileExists{microtype.sty}{% use microtype if available
  \usepackage[]{microtype}
  \UseMicrotypeSet[protrusion]{basicmath} % disable protrusion for tt fonts
}{}
% Make \paragraph and \subparagraph free-standing
\makeatletter
\ifx\paragraph\undefined\else
  \let\oldparagraph\paragraph
  \renewcommand{\paragraph}{
    \@ifstar
      \xxxParagraphStar
      \xxxParagraphNoStar
  }
  \newcommand{\xxxParagraphStar}[1]{\oldparagraph*{#1}\mbox{}}
  \newcommand{\xxxParagraphNoStar}[1]{\oldparagraph{#1}\mbox{}}
\fi
\ifx\subparagraph\undefined\else
  \let\oldsubparagraph\subparagraph
  \renewcommand{\subparagraph}{
    \@ifstar
      \xxxSubParagraphStar
      \xxxSubParagraphNoStar
  }
  \newcommand{\xxxSubParagraphStar}[1]{\oldsubparagraph*{#1}\mbox{}}
  \newcommand{\xxxSubParagraphNoStar}[1]{\oldsubparagraph{#1}\mbox{}}
\fi
\makeatother


\usepackage{longtable,booktabs,array}
\usepackage{calc} % for calculating minipage widths
% Correct order of tables after \paragraph or \subparagraph
\usepackage{etoolbox}
\makeatletter
\patchcmd\longtable{\par}{\if@noskipsec\mbox{}\fi\par}{}{}
\makeatother
% Allow footnotes in longtable head/foot
\IfFileExists{footnotehyper.sty}{\usepackage{footnotehyper}}{\usepackage{footnote}}
\makesavenoteenv{longtable}
\usepackage{graphicx}
\makeatletter
\newsavebox\pandoc@box
\newcommand*\pandocbounded[1]{% scales image to fit in text height/width
  \sbox\pandoc@box{#1}%
  \Gscale@div\@tempa{\textheight}{\dimexpr\ht\pandoc@box+\dp\pandoc@box\relax}%
  \Gscale@div\@tempb{\linewidth}{\wd\pandoc@box}%
  \ifdim\@tempb\p@<\@tempa\p@\let\@tempa\@tempb\fi% select the smaller of both
  \ifdim\@tempa\p@<\p@\scalebox{\@tempa}{\usebox\pandoc@box}%
  \else\usebox{\pandoc@box}%
  \fi%
}
% Set default figure placement to htbp
\def\fps@figure{htbp}
\makeatother


% definitions for citeproc citations
\NewDocumentCommand\citeproctext{}{}
\NewDocumentCommand\citeproc{mm}{%
  \begingroup\def\citeproctext{#2}\cite{#1}\endgroup}
\makeatletter
 % allow citations to break across lines
 \let\@cite@ofmt\@firstofone
 % avoid brackets around text for \cite:
 \def\@biblabel#1{}
 \def\@cite#1#2{{#1\if@tempswa , #2\fi}}
\makeatother
\newlength{\cslhangindent}
\setlength{\cslhangindent}{1.5em}
\newlength{\csllabelwidth}
\setlength{\csllabelwidth}{3em}
\newenvironment{CSLReferences}[2] % #1 hanging-indent, #2 entry-spacing
 {\begin{list}{}{%
  \setlength{\itemindent}{0pt}
  \setlength{\leftmargin}{0pt}
  \setlength{\parsep}{0pt}
  % turn on hanging indent if param 1 is 1
  \ifodd #1
   \setlength{\leftmargin}{\cslhangindent}
   \setlength{\itemindent}{-1\cslhangindent}
  \fi
  % set entry spacing
  \setlength{\itemsep}{#2\baselineskip}}}
 {\end{list}}
\usepackage{calc}
\newcommand{\CSLBlock}[1]{\hfill\break\parbox[t]{\linewidth}{\strut\ignorespaces#1\strut}}
\newcommand{\CSLLeftMargin}[1]{\parbox[t]{\csllabelwidth}{\strut#1\strut}}
\newcommand{\CSLRightInline}[1]{\parbox[t]{\linewidth - \csllabelwidth}{\strut#1\strut}}
\newcommand{\CSLIndent}[1]{\hspace{\cslhangindent}#1}



\setlength{\emergencystretch}{3em} % prevent overfull lines

\providecommand{\tightlist}{%
  \setlength{\itemsep}{0pt}\setlength{\parskip}{0pt}}



 


\usepackage{graphicx}
\usepackage{bbm}
\usepackage{amssymb}
\usepackage{amsmath}
\usepackage[nomarginpar]{geometry}
\usepackage{setspace}
%\usepackage{natbib}
\usepackage{lscape}
\usepackage{rotating}
\usepackage{tabularx, calc}
\usepackage{threeparttable}
\usepackage{picinpar}
\usepackage{longtable}
\usepackage{ae}
\usepackage{float}
\usepackage{morefloats}
\usepackage{palatino}
\usepackage{hyperref}
\usepackage{color}
\usepackage{adjustbox}
\usepackage{booktabs}
\usepackage[normalem]{ulem}
\usepackage{ragged2e}
\usepackage{changepage}
\usepackage{longtable}
%\usepackage{graphics}
%\usepackage[tableposition=top]{caption}

\geometry{left=1.0in,right=1.0in,top=1.0in,bottom=1.0in}
\onehalfspacing
\usepackage{float}
\usepackage{tabularray}
\usepackage[normalem]{ulem}
\usepackage{graphicx}
\usepackage{rotating}
\UseTblrLibrary{siunitx}
\NewTableCommand{\tinytableDefineColor}[3]{\definecolor{#1}{#2}{#3}}
\newcommand{\tinytableTabularrayUnderline}[1]{\underline{#1}}
\newcommand{\tinytableTabularrayStrikeout}[1]{\sout{#1}}
\makeatletter
\@ifpackageloaded{tcolorbox}{}{\usepackage[skins,breakable]{tcolorbox}}
\@ifpackageloaded{fontawesome5}{}{\usepackage{fontawesome5}}
\definecolor{quarto-callout-color}{HTML}{909090}
\definecolor{quarto-callout-note-color}{HTML}{0758E5}
\definecolor{quarto-callout-important-color}{HTML}{CC1914}
\definecolor{quarto-callout-warning-color}{HTML}{EB9113}
\definecolor{quarto-callout-tip-color}{HTML}{00A047}
\definecolor{quarto-callout-caution-color}{HTML}{FC5300}
\definecolor{quarto-callout-color-frame}{HTML}{acacac}
\definecolor{quarto-callout-note-color-frame}{HTML}{4582ec}
\definecolor{quarto-callout-important-color-frame}{HTML}{d9534f}
\definecolor{quarto-callout-warning-color-frame}{HTML}{f0ad4e}
\definecolor{quarto-callout-tip-color-frame}{HTML}{02b875}
\definecolor{quarto-callout-caution-color-frame}{HTML}{fd7e14}
\makeatother
\makeatletter
\@ifpackageloaded{caption}{}{\usepackage{caption}}
\AtBeginDocument{%
\ifdefined\contentsname
  \renewcommand*\contentsname{Table of contents}
\else
  \newcommand\contentsname{Table of contents}
\fi
\ifdefined\listfigurename
  \renewcommand*\listfigurename{List of Figures}
\else
  \newcommand\listfigurename{List of Figures}
\fi
\ifdefined\listtablename
  \renewcommand*\listtablename{List of Tables}
\else
  \newcommand\listtablename{List of Tables}
\fi
\ifdefined\figurename
  \renewcommand*\figurename{Figure}
\else
  \newcommand\figurename{Figure}
\fi
\ifdefined\tablename
  \renewcommand*\tablename{Table}
\else
  \newcommand\tablename{Table}
\fi
}
\@ifpackageloaded{float}{}{\usepackage{float}}
\floatstyle{ruled}
\@ifundefined{c@chapter}{\newfloat{codelisting}{h}{lop}}{\newfloat{codelisting}{h}{lop}[chapter]}
\floatname{codelisting}{Listing}
\newcommand*\listoflistings{\listof{codelisting}{List of Listings}}
\makeatother
\makeatletter
\makeatother
\makeatletter
\@ifpackageloaded{caption}{}{\usepackage{caption}}
\@ifpackageloaded{subcaption}{}{\usepackage{subcaption}}
\makeatother
\usepackage{bookmark}
\IfFileExists{xurl.sty}{\usepackage{xurl}}{} % add URL line breaks if available
\urlstyle{same}
\hypersetup{
  pdftitle={More money, more effect? Employment effects of job search services of different intensity},
  pdfauthor={Álvaro F. Junquera},
  pdfkeywords={active labour market policies, job search assistance, job
search intermediation, public employment services, voucher},
  colorlinks=true,
  linkcolor={blue},
  filecolor={Maroon},
  citecolor={blue},
  urlcolor={blue},
  pdfcreator={LaTeX via pandoc}}


\title{More money, more effect? Employment effects of job search
services of different intensity}
\author{Álvaro F. Junquera}
\date{March 30, 2026}
\begin{document}
\def\spacingset#1{\renewcommand{\baselinestretch}%
{#1}\small\normalsize} \spacingset{1}


%%%%%%%%%%%%%%%%%%%%%%%%%%%%%%%%%%%%%%%%%%%%%%%%%%%%%%%%%%%%%%%%%%%%%%%%%%%%%%

\date{\href{https://osf.io/preprints/socarxiv/rjshu}{Link to Latest version}\\ \vspace{1em}  Last updated March
30, 2026}
\title{More money, more effect? Employment effects of job search
services of different intensity}
\author{
Álvaro F. Junquera\thanks{Correspondence: alvaro.junquera@bsc.es, ORCID:
0000-0003-0080-7608. We thank the collaboration of the public employment
office of Veneto. We also appreciate the suggestions offered by Enrico
Rettore, the help offered by Jungmo Yoon, and the commentaries made by
discussants of the 1st Workshop on Comparative Social and Public Policy,
the 16th Spanish Labour Economics meeting, the 39th Conference of the
Italian Association of Labour Economics, and the 4th Workshop on
International Trends in Economic Research. This article has been
developed as part of the PhD programme in Sociology of the Universitat
Autònoma de Barcelona (UAB) and the PhD programme in Economics and
Management of the Università degli Studi di Padova (UNIPD).}\\
Barcelona Supercomputing Center (BSC-CNS)\\
}
\maketitle

\bigskip
\bigskip
 % added by AFJ
\begin{abstract}
Job search services modestly raise employment, but it is far from clear
whether spending more on each participant improves outcomes further. We
address this question by studying Assegno per il Lavoro, an Italian
active labour market programme that assigns jobseekers to treatment
groups with increasing endowments of job search services. Each
intervention is made up of a voucher to fund job search assistance and a
performance-based payment related to job search intermediation.
Participants are assigned to groups with treatment endowments that
increase according to a scoring variable. Leveraging this design, we
apply a regression discontinuity analysis to estimate effects on both
employment duration and employment quality. We find that providing a
larger voucher for job search assistance does not increase take-up of
the service, limiting its potential impact on employment. On employment
duration, effects of higher endowments after 18 months are essentially
null for both the average and the quartiles of days employed. On
employment quality, we find that the intermediate endowment slightly
increases the length of the contracts. Moreover, a design with different
performance payments for job search intermediation seems to deter cream
skimming, since workers are not discriminated in their access to the
services. These results highlight both the risks of vouchers to
implement public interventions and the potential of differentiated
performance payments to deter cream skimming.
\end{abstract}
 % added by AFJ

\bigskip
\noindent%
{\it Keywords:} active labour market policies, job search
assistance, job search intermediation, public employment
services, voucher
\vfill

 % added by AFJ
\noindent%
{\it JEL codes:} J64, J68, C21
\vfill


\newpage
\spacingset{1.2} % DON'T change the spacing!

\section{Introduction}\label{sec-intro}

Job search services are a type of active labour market policy (ALMP)
usually qualified as effective to slightly increase the probability of
employment (\citeproc{ref-cardetal2018}{Card et al. 2018};
\citeproc{ref-voorenetal2019}{Vooren et al. 2019}). However, designing a
programme like this requires knowledge of the specific tools and design
features that optimize the average treatment effect. On the tools,
recent research proposes to separate job search assistance from
intermediation, suggesting that they influence the outcome through
different mechanisms (\citeproc{ref-cheungetal2025}{Cheung et al.
2025}). This distinction has not been addressed in previous
meta-analyses, which have overlooked the differences between the
assistance and intermediation mechanisms. In contrast, we analyze a
programme that allows a particular focus on job search intermediation, a
mechanism underexplored in the prior literature. On the design features,
aspects like the intensity of the actions have been discussed when
comparing different programmes. Literature reviews in economics
(\citeproc{ref-voorenetal2019}{Vooren et al. 2019}) and psychology
(\citeproc{ref-liuetal2014}{Liu et al. 2014}) found that, in highly
developed countries, the mean duration of an ALMP was not related to its
employment effect. However, Escudero et al.
(\citeproc{ref-escuderoetal2019}{2019}) detected the opposite for low
and middle-income countries. This finding might suggest an interaction
between the intensity of the intervention and other characteristics of
the labour markets. If intensity effects are indeed null in developed
countries, this would raise the question on the minimum treatment dose
needed to generate a certain effect. Evidence of null intensity effects
would justify reallocating resources toward covering more jobseekers
rather than intensifying treatment for each participant.

This article studies the employment effects of different treatments with
diverse doses of job search services. We evaluate an active labour
market programme called Assegno per il Lavoro and implemented in Veneto
(Italy) that offers different treatment endowments. Each participant of
the treatment group is linked to a voucher of job search assistance and
to services of job search intermediation. Importantly, both the voucher
value and the intermediation payment increase across groups. Concerning
the duration of employment, we study the effects of different treatment
endowments on the mean and different quantiles of the distribution of
days worked. Regarding employment quality, we study the effects on the
probability of signing different types of employment contracts. Thanks
to the assignment mechanism of the programme, we carry out a sharp
regression discontinuity analysis that allows us to estimate local
effects at the threshold. Identification was possible for two treatment
comparisons: intermediate (\(B\)) vs.~low endowment (\(A\)) and high
(\(C\)) vs.~intermediate endowment (\(B\)).

The results on the employment effects show that increasing the endowment
is in general ineffective to increase the number of days worked, but it
avoids the discrimination of hard-to-place jobseekers. In terms of
employment duration, the impact on the average of days worked is
eventually null at all semesters for the first treatment comparison
(\(B\) vs.~\(A\)) and of 10 extra days during the first year for the
second treatment comparison (\(C\) vs.~\(B\)). Looking at the median of
the distributions, we estimate that there is a gain of employment in the
intermediate-low contrast during the third semester and in the
high-intermediate comparison during the second semester. The drawback is
that all these effects vanish subsequently. In terms of employment
quality, the intermediate-low comparison triggers a higher probability
of getting an open-ended contract during the implementation period of
the programme. In contrast, this is not detected for the
high-intermediate comparison, which generates an increase in the
probability of signing a short temporary contract as the longest
employment episode. Concerning programme implementation, on the one
hand, we detect a very small change in the actual duration of job search
assistance at the threshold, consistent with prior evidence on
implementation failures in voucher-based public services. Consequently,
job search intermediation appears to be responsible for the few impacts
observed. On the other hand, under the differentiated payment structure,
we find no evidence of discrimination in access to job search services
for the less-employable jobseekers.

\textbf{Related literature}. There is a broad literature estimating the
effect of participating in job search services on employment attainment
both in economics (\citeproc{ref-cardetal2018}{Card et al. 2018};
\citeproc{ref-voorenetal2019}{Vooren et al. 2019}) and psychology
(\citeproc{ref-liuetal2014}{Liu et al. 2014}). Our work contributes to
three research lines that deal with the employment effects of the
following design features: the intensity of assistance, the intensity of
intermediation, and the private provision of public employment services.

Firstly, job search assistance (JSA) is aimed at improving the job
search of the unemployed individual by intervening on her job search
skills, motivation, or targets. The heterogeneous effects of JSA in
function of its intensity have been rarely studied in an European
context. Liu et al. (\citeproc{ref-liuetal2014}{2014}) performed a
meta-analysis on the employment effects of JSA including 32 randomized
experiments and 15 observational studies, most of them based on the USA.
They concluded that there was not an association between the intensity
(duration) of the intervention and the effect magnitude. Martinson et
al. (\citeproc{ref-martinsonetal2019}{2019}) and Dyke et al.
(\citeproc{ref-dykeetal2006}{2006}) also found a null effect of JSA
intensity on employment with studies almost exclusively from the USA. On
the contrary, Meyer (\citeproc{ref-meyer1995}{1995}) detected a positive
employment effect also in the USA. More research is needed for the
European countries, as there are substantial differences in the
organization of public administrations and the markets for ALMPs
providers between the USA and Europe.

Secondly, job search intermediation (JSI) involves a delegation of the
job search from the worker to a third party, whether as a full
delegation or as a complementary service. There are only a few
evaluations of these interventions, but they usually conclude positive
effects on employment (\citeproc{ref-fougereetal2009}{Fougère et al.
2009}; \citeproc{ref-engstrometal2012}{Engström et al. 2012};
\citeproc{ref-vandenbergetal2019}{Van Den Berg et al. 2019};
\citeproc{ref-cheungetal2025}{Cheung et al. 2025}). Three other articles
have investigated the effect of JSI under performance-based payments.
Zanella and Salomone (\citeproc{ref-zanellasalomone2025}{2025}) found a
positive impact in private providers, whereas Koning and Heinrich
(\citeproc{ref-koningheinrich2013}{2013}) concluded the same under a mix
of ALMPs but only for non-disable jobseekers. Burgess et al.
(\citeproc{ref-burgessetal2017}{2017}) studied public providers and only
detected positive effects in large establishments.

Thirdly, researchers have also studied the potential benefits and risks
of a private provision of public employment services via vouchers and
performance pays (\citeproc{ref-hippwarner2008}{Hipp and Warner 2008}).
Regarding a provision with vouchers, in the American context, Barnow
(\citeproc{ref-barnow2009}{2009}) reviewed evaluations of voucher
implementations in adult training, noting mixed estimates on the
employment effects. Doerr et al. (\citeproc{ref-doerretal2017}{2017})
found that, on average, receiving a voucher generated large and lasting
negative effects on employment. Like Huber et al.
(\citeproc{ref-huberetal2018}{2018}), they found small positive effects
emerging in the very long run. In the European context, Brenninkmeijer
and Blonk (\citeproc{ref-brenninkblonk2012}{2012}) conducted an
experiment on JSA in the Netherlands, evaluating an open-dose (voucher)
against a closed-dose treatment and a no-treatment group. In the short
term, the voucher condition yielded results similar to the control
condition. An Italian programme evaluated by Linfante et al.
(\citeproc{ref-linfanteetal2019}{2019}) found null effects on employment
one year after treatment for compliers who registered after receiving an
invitation.

Concerning a provision with performance pays, one of the potential
disadvantages of a quasi-marketization of public services is the cream
skimming problem. Cream skimming refers to the rejection of
hard-to-place jobseekers by service providers. The evidence on the
existence of this problem is mixed, with some studies finding it for
certain groups (\citeproc{ref-koningheinrich2013}{Koning and Heinrich
2013}), other studies rejecting its existence
(\citeproc{ref-aakviketal2005}{Aakvik et al. 2005}) and other articles
highlighting the measurement problems of this concept
(\citeproc{ref-heckmansmith2004}{Heckman and Smith 2004}). Establishing
different performance pays in function of the jobseeker's employability
has been sometimes presented as a shelter against this problem.

This paper makes three contributions to the literature. First, we
provide new evidence on the dose--response function in job search
assistance within a European context, where programme designs and labour
market institutions differ substantially from the US-based prior
literature. Second, we study the heterogeneous effects of job search
intermediation under performance-based payments, focusing for the first
time on how the intensity of such payments shapes employment outcomes.
Third, we assess whether differentiated performance payments succeed in
discouraging cream skimming, contributing to the literature on the
design of quasi-markets for public employment services. On the
methodological side, our analysis uses a credible identification that
combines regression discontinuity with local quantile treatment effects,
one of the first applications of these recently proposed techniques.

We continue in Section~\ref{sec-policy} describing the programme design
evaluated. In Section~\ref{sec-meth}, we present the data and the
methods employed to answer our research questions.
Section~\ref{sec-results} shows the results of our work and assesses the
plausibility of our identification assumptions. Lastly, we draw some
conclusions in Section~\ref{sec-conc}.

\section{Policy design}\label{sec-policy}

Assegno per il Lavoro (hereinafter, AXL) was a programme developed in
Veneto (Italy) from July 2017 to May 2022 and funded by the public
employment office of the region. The general objective of AXL was to
reduce unemployment. Initially, the target group were those over 35
years old that met at least one of the conditions of social
vulnerability included in Government of Veneto
(\citeproc{ref-veneto2017a}{2017a}). From April 2018, the programme
extended to all unemployed individuals over 30 years old.

To achieve this objective, policymakers designed an intervention with
diverse tools or components. The individual could receive job search
assistance, adult training, job search intermediation, and internships.
The intensity of such intervention could vary according to the treatment
group. Table~\ref{tbl-endowdur} shows the amount of money available for
each participant related to each component by treatment group. There are
three main sources of variation between groups. Regarding job search
assistance, group \(B\) almost doubles the number of maximum hours
available compared to group \(A\). We can see a remarkable difference on
this matter between group \(C\) and group \(B\) as well. Concerning job
search intermediation and internships, the money devoted to such
components also changes by treatment group. Lastly, the maximum duration
of the intervention also increases from \(A\) to \(C\), although the
differences are not substantial.

\begin{longtable}[]{@{}
  >{\raggedright\arraybackslash}p{(\linewidth - 6\tabcolsep) * \real{0.2703}}
  >{\raggedright\arraybackslash}p{(\linewidth - 6\tabcolsep) * \real{0.2432}}
  >{\raggedright\arraybackslash}p{(\linewidth - 6\tabcolsep) * \real{0.2432}}
  >{\raggedright\arraybackslash}p{(\linewidth - 6\tabcolsep) * \real{0.2432}}@{}}
\caption{Endowments and duration of AXL by treatment group
(\citeproc{ref-veneto2017a}{Government of Veneto 2017a};
\citeproc{ref-veneto2018}{Government of Veneto
2018})}\label{tbl-endowdur}\tabularnewline
\toprule\noalign{}
\begin{minipage}[b]{\linewidth}\raggedright
\end{minipage} & \begin{minipage}[b]{\linewidth}\raggedright
Group A
\end{minipage} & \begin{minipage}[b]{\linewidth}\raggedright
Group B
\end{minipage} & \begin{minipage}[b]{\linewidth}\raggedright
Group C
\end{minipage} \\
\midrule\noalign{}
\endfirsthead
\toprule\noalign{}
\begin{minipage}[b]{\linewidth}\raggedright
\end{minipage} & \begin{minipage}[b]{\linewidth}\raggedright
Group A
\end{minipage} & \begin{minipage}[b]{\linewidth}\raggedright
Group B
\end{minipage} & \begin{minipage}[b]{\linewidth}\raggedright
Group C
\end{minipage} \\
\midrule\noalign{}
\endhead
\bottomrule\noalign{}
\endlastfoot
Total endowment & 3,536 € & 4,264 € & 5,796 € \\
- Job search assistance & 266 € (7 hours) & 494 € (13 hours) & 1,026 €
(27 hours) \\
- Adult training & 1,770 € & 1,770 € & 1,770 € \\
- Job search intermediation* & 1,500 € & 2,000 € & 3,000 € (for JSI and
internships) \\
- Internships* & 0 € & 0 € & \\
Maximum duration & 3 months (6 with interruption) & 5 months (8 with
interruption) & 6 months (9 with interruption) \\
\end{longtable}

The implementation process is carried out through vouchers and
performance-based pays. The individual decides first where to spend it.
There is a quasi-market of private providers that offer all these
services. The public employment office lists all of them and the
individual chooses where to go. On the voucher side, participant and
provider decide then how much to spend on job search assistance and on
adult training. The hourly price of the first component is fixed, but it
is not for the second one. On the prize side, job search intermediation
and internships are not paid depending on the service provision. It is
strictly a performance-based payment: the provider only gets the money
if the participant signs an employment contract. Moreover, the
performance-based pay increases in function of the duration of the
contract\footnote{Table~\ref{tbl-endowdur} shows the performance-based
  pay for JSI in case of an open-ended contract. If a fixed-term
  contract of more than 12 months is signed, this quantity is multiplied
  by 0.8. If the duration is between 6 and 12 months, the factor is 0.6.}.

Table~\ref{tbl-endowdur} shows that the endowment of the programme
depends on the treatment group in which the participant is allocated. In
AXL, the individual is assigned to a certain treatment group according
to a scoring variable that estimates the probability of being employed
in 24 months. The score has the following functional form
(\citeproc{ref-veneto2017b}{Government of Veneto 2017b}):

\[
\begin{aligned}
\hat{S} i &= 1/[1 + \exp(-(\hat{\beta}_0 + \hat{\beta}_1 sex_i + \hat{\beta}_2 age_i + \hat{\beta}_3 foreign_i + \hat{\beta}_4 studies_i + \hat{\beta}_5 disab_i \\
&\quad + \hat{\beta}_6 termination_i + \hat{\beta}_7 occup_i + \hat{\beta}_8 sector_i))]
\end{aligned}
\]

\(\hat{S}_i\) incorporates both sociodemographic covariates (sex, age,
being foreign, level of studies, and having a disability) and covariates
on the last employment (reason of termination, major group of
occupation, and sector). To make the scoring variable following the same
direction of treatment endowment, we reversed the original scoring
variable generating \(S_i\). Once the participant receives a score, it
is allocated to a certain treatment group according to the next
function:

\[
group_i =
\begin{cases}
A & 0 \le s \le 0.4 \\
B & 0.41 \le s \le 0.6 \\
C & 0.61 \le s \le 1
\end{cases}
\]

\section{Methods}\label{sec-meth}

\subsection{Data}\label{data}

The data comes from administrative records of Veneto Lavoro, the public
employment service of the region of Veneto (Italy). We use microdata on
all individuals that applied to AXL during the whole implementation
period of the programme.

We have information on three outcome variables measured at different
periods. Firstly, related to the employment duration, we know the number
of days worked during the first, the second, the third, and the fourth
semester since the treatment started. Secondly, regarding employment
quality, we have a polytomous outcome variable that registers what was
the longest employment contract signed since the treatment began. Notice
that the type of contract is considered one of the two dimensions of
employment quality by Muñoz De Bustillo et al.
(\citeproc{ref-munozdebustilloetal2011}{2011}). We distinguished five
categories: open-ended contract (OEC), temporary contract of long
duration (i.e.~more than 12 months), temporary contract of intermediate
duration (i.e.~from 6 to 12 months), temporary contract of short
duration (5 or less months), and no employment. We also have information
about the type of contract linked to the performance-base pay of
intermediation. In some cases, such contract may be the same as the
longest employment contract, but not necessarily. Thirdly, related to
treatment uptake, we know the doses of treatment actually received by
participants. We have the number of hours of job search assistance and
the number of hours of adult training. This last set of outcomes will
help us disentangle some possible mechanisms to explain the results.

The remaining covariates in our dataset are the scoring variable and
other sociodemographic variables. The scoring or running variable is the
unique confounder according to the programme design, so we do not need
additional control variables for identification. The remaining
sociodemographic covariates will be employed for the plausibility and
heterogeneity analyses. In the supplementary document, we include the
full list of variables with their definitions.

\subsection{Analysis}\label{analysis}

\subsubsection{Target parameters}\label{target-parameters}

We consider different outcomes, so we denote as \(Y_{oi}\) the response
variable for outcome \(o\) and individual \(i\). For a simpler
exposition, we remove the subindex \(o\) without loss of generality.
During the presentation of results, we also eliminate subindex \(i\) to
facilitate the reading. We denote the potential outcome for the state of
the world with the active treatment as \(Y_i(1)\) and the potential
outcome for the state of the world without that active treatment as
\(Y_i(0)\).

We are interested in two treatment comparisons. \(D_{1i}\) is a binary
variable that collects if individual was treated under group \(B\) (with
value 1) or under group \(A\) (with value 0). \(D_{2i}\) is another
binary variable that informs if individual was treated under group \(C\)
(with value 1) or under group \(B\) (with value 0). For each case,
treatment with value 1 is referred as the active treatment, whereas
treatment with value 0 is called the control treatment. As we mentioned
before, treatment group is determined by the scoring variable \(S_i\).
For each treatment comparison, we partition the sample by delimiting the
range of the score and centring it at the cutoff, so both \(S_{1i}\)
(affecting \(D_{1i}\)) and \(S_{2i}\) (affecting \(D_{2i}\)) present the
cutoff at \(c=0\). Therefore, we may denote more generally \(D_{zi}\)
and \(S_{zi}\) for \(z \in \{1,2\}\), but in the following we omit
subindex \(z\) to reduce the notational burden.

Our target parameters are the local average treatment effect of \(D_i\)
on \(Y_i\) (\(LATE\)), the local average treatment effect of \(D_i\) on
the probability of \(Y_i=j\) (\(LATE_j\)), and the local quantile
treatment effect for a certain quantile \(\mathbb{Q}_\tau\) at
percentile \(\tau\) of \(D_i\) on \(Y_i\) (\(LQTE(\tau)\)). In all
cases, by local we mean the subpopulation with a score just at the
threshold, which determines receiving the higher endowment in each
comparison. Equivalently, these target parameters can be understood as
local effects of a certain endowment or as intention-to-treat effects of
a certain treatment. Formally, we can write the estimands as

\[
LATE = \mathbb{E}[Y_i(1)|S_{zi}=c] - \mathbb{E}[Y_i(0)|S_{zi}=c] \text{ for } z \in \{1, 2\}
\] \[
LATE_j = \mathbb{P}[Y_i(1)=j|S_{zi}=c] - \mathbb{P}[Y_i(0)=j|S_{zi}=c] \text{ for } z \in \{1, 2\}
\] \[
LQTE(\tau) = \mathbb{Q}_{\tau}[Y_i(1)|S_{zi}=c] - \mathbb{Q}_{\tau}[Y_i(0)|S_{zi}=c] \text{ for } z \in \{1, 2\}
\]

\subsubsection{Identification}\label{identification}

Since the target parameters contain unobservable potential outcomes, we
require identification assumptions to estimate them using observable
potential outcomes. Following Rubin (\citeproc{ref-rubin2007}{2007}), an
identification strategy is more credible when it is closer to the
properties of a randomized experiment. Therefore, we propose to leverage
the design of AXL to use a sharp regression discontinuity strategy. This
method is based on two identification assumptions related to two
properties of randomized experiments.

\begin{tcolorbox}[enhanced jigsaw, breakable, leftrule=.75mm, arc=.35mm, left=2mm, colback=white, rightrule=.15mm, colframe=quarto-callout-tip-color-frame, toprule=.15mm, bottomrule=.15mm, opacityback=0]

\vspace{-3mm}\textbf{Assumption 1}\vspace{3mm}

The assignment mechanism is \(D_i = \mathbf{1}(S_i \ge c)\), so
\(\lim_{s \downarrow c} \mathbb{E}[D_i|S_i=s] - \lim_{s \uparrow c} \mathbb{E}[D_i|S_i=s] \ne 0\).

\end{tcolorbox}

The first identification assumption is the unconfoundedness of the
assignment mechanism, which departs from the knowledge on how the
treatments were assigned (\citeproc{ref-abadiecattaneo2018}{Abadie and
Cattaneo 2018}; \citeproc{ref-frolichsperlich2019}{Frölich and Sperlich
2019}). The main idea is that adjusting for the scoring variable, the
assignment variable is not related to any potential outcome. In our
case, it is a weak assumption thanks to the detailed knowledge on the
treatment assignment function. Both the awareness of the assignment
function and the unconfoundedness of the assignment mechanism are
properties of randomized experiments.

If treatment assignment was not only influenced by the score, but also
by other unadjusted variables like ability, Assumption 1 would be false.
This problem is known as score manipulation
(\citeproc{ref-frolichsperlich2019}{Frölich and Sperlich 2019}), so we
will examine how feasible is to manipulate the specific variables that
determine the score used in AXL.

\begin{tcolorbox}[enhanced jigsaw, breakable, leftrule=.75mm, arc=.35mm, left=2mm, colback=white, rightrule=.15mm, colframe=quarto-callout-tip-color-frame, toprule=.15mm, bottomrule=.15mm, opacityback=0]

\vspace{-3mm}\textbf{Assumption 2}\vspace{3mm}

\begin{itemize}
\tightlist
\item
  For \(LATE\): \(\mathbb{E}[Y_i(1)|S_i]\) is continuous at \(S_i = c\)
  and \(\mathbb{E}[Y_i(0)|S_i]\) is continuous at \(S_i = c\).
\item
  For \(LATE_j\): \(\mathbb{P}[Y_i(1)=j|S_i]\) is continuous at
  \(S_i = c\) and \(\mathbb{P}[Y_i(0)=j|S_i]\) is continuous at
  \(S_i = c\).
\item
  For \(LQTE(\tau)\): \(F_{Y(1)|S}(y|c)\) is continuous and
  \(F_{Y(0)|S}(y|c)\) is continuous.
\end{itemize}

\end{tcolorbox}

The second identification assumption solves the absence of overlap
through the continuity of potential outcome functions
(\citeproc{ref-frolichsperlich2019}{Frölich and Sperlich 2019}). While
the designs of randomized experiments present probabilistic assignment
mechanisms, the designs linked to the sharp regression discontinuity are
based on a deterministic assignment mechanism. There are two approaches
to bypass this problem in the regression discontinuity context
(\citeproc{ref-cattaneoetal2019}{Cattaneo et al. 2019}). The local
randomization approach is suitable when we have a natural experiment as
defined by Titiunik (\citeproc{ref-titiunik2021}{2021}). It is not our
case, because our assignment mechanism does not incorporate any element
of chance. The continuity approach to regression discontinuity solves
the lack of overlap by performing a small extrapolation. The rationale
is that extrapolations distant from the observed data are more
model-dependent (\citeproc{ref-kingzen2006}{King and Zeng 2006}), so
this strategy focuses on a more modest estimand around the cutoff.

There are two situations that might provoke violations of Assumption 2.
On the one hand, if another programme employed the same scoring variable
to allocate treatments and its effect was not null. Thereby,
\(\mathbb{E}[Y_i(0)|S_i]\) or its homologous would suffer a
discontinuity at the cutoff, because this would be the point at which
the alternative treatment is activated. Two checks will be performed to
assess the plausibility of this situation. First, we will use
pseudo-treatment or placebo outcome tests just at the cutoff
(\citeproc{ref-frolichsperlich2019}{Frölich and Sperlich 2019};
\citeproc{ref-eggersetal2024}{Eggers et al. 2024}). Following Manski and
Pepper (\citeproc{ref-manskipepper2018}{2018}), a statistically
significant difference between groups would diminish the credibility of
the continuity assumption. Second, we will examine the documents of the
AXL design to judge the possibility of other programmes employing the
same running variable. On the other hand, if treatment groups in our
sample had different conditional expected values in covariates included
in the confounding variable, the credibility of the continuity
assumption would also be lower. Therefore, we will study the balance of
these covariates among the different treatment groups.

\subsubsection{Estimation and inference}\label{estimation-and-inference}

We are going to estimate the target parameters using graphical and
quantitative methods. As a first step, we use regression discontinuity
plots both with local mean/proportions and with local medians
(\citeproc{ref-cattaneoetal2024}{Cattaneo et al. 2024}). As a second
step, we analyse the data with more fine-grained estimation and
inference techniques.

Regarding point estimation, we have used three estimators, one for each
estimand. For quantitative outcomes, we can credibly estimate both the
\(LATE\) and the \(LQTE(\tau)\). In order to estimate the \(LATE\), we
have used the local polynomial estimator Cattaneo and Titiunik
(\citeproc{ref-cattaneotitiunik2022}{2022}). A common bandwidth for both
sides has been chosen by minimizing the mean squared error following
standard procedures (\citeproc{ref-calonicoetal2017}{Calonico et al.
2017}). For each \(LATE\), two specifications have been estimated, one
quadratic and other linear in the score, both with the triangular
kernel\footnote{We chose the triangular kernel because of its point
  estimation optimality property, its consistency in spirit with the
  available kernels of Qu and Yoon (\citeproc{ref-quyoon2019}{2019}) and
  its compatibility with the procedure of Pei et al.
  (\citeproc{ref-peietal2022}{2022}). The results are essentially the
  same with a uniform kernel for almost all semesters.}. The definitive
specification has been selected according to the asymptotic mean squared
error as recommended by Pei et al. (\citeproc{ref-peietal2022}{2022}).

To estimate the \(LQTE(\tau)\), we have applied the local estimator
proposed by Qu and Yoon (\citeproc{ref-quyoon2019}{2019}). The procedure
to choose the bandwidth is adapted to the quantile setting but preserves
the same intuition: the minimization of the mean squared error. We have
estimated the optimal bandwidths according to three procedures proposed
by Qu and Yoon (\citeproc{ref-quyoon2015}{2015}) and Qu and Yoon
(\citeproc{ref-quyoon2019}{2019}): treating the cutoff as an interior
point, as a boundary point, and with the quantile adaptation of the
bandwidth selector by Imbens and Kalyanaraman
(\citeproc{ref-imbenskaly2012}{2012}). Then, to avoid bias, we have
selected the lowest bandwidth among the three alternatives. As
recommended by the authors of the method, all specifications consider a
linear specification of the regression coefficients of \(S_i\) and
employ the Epanechnikov kernel. For quantile regressions, we do not
maintain the rank invariance assumption, so interpretation of results
must be made in accordance with that. Under rank dissimilarity,
estimates of \(LQTE(\tau)\) are not informative about on which
individual the treatment is most effective
(\citeproc{ref-dongshen2018}{Dong and Shen 2018}). Nonetheless, they
inform about changes at the potential outcome distributions, which helps
to assess the usefulness of the programme as a whole.

For qualitative outcomes, we can only estimate the \(LATE_j\). We have
employed the local logit model with maximum likelihood estimator
developed by Xu (\citeproc{ref-xu2017}{2017}). The bandwidth has also
been chosen by minimizing the mean squared error, with a procedure
adapted to the categorical setting by Xu (\citeproc{ref-xu2017}{2017}).
His method also departs from a linear specification of the coefficients
of \(S_i\) and uses a uniform kernel.

Concerning inference, methodologists have considered that non-parametric
estimators minimize bias \emph{and} variance when optimizing the mean
squared error, so some bias remains. Following Cattaneo et al.
(\citeproc{ref-cattaneoetal2019}{2019}), we use a bandwidth that is
optimal for point estimation but removing the remaining bias for
inference aims. For the mean estimand, we executed a Wald test of no
effect employing the robust bias-corrected approach of Calonico et al.
(\citeproc{ref-calonicoetal2014}{2014}). For the probability estimand,
Xu (\citeproc{ref-xu2017}{2017}) adapts the approach of Calonico et al.
(\citeproc{ref-calonicoetal2014}{2014}) to the categorical setting also
with a Wald test. For the quantile estimand, the first inference
procedure contrasts treatment non-significance and is a Wald test of
\(H_0:LQTE(\tau)=0\) for all \(\tau \in \mathcal{T}\), with
\(\mathcal{T}\) denoting the quantile range, and
\(H_1:LQTE(\tau) \ne 0\) for some \(\tau \in \mathcal{T}\). The second
inference procedure contrasts the homogeneity of treatment effects
across quantiles and is a Wald test of \(H_0:LQTE(\tau)\) is constant
over \(\tau \in \mathcal{T}\), with \(H_1:LQTE(\tau) \ne LQTE(v)\) for
some \(\tau,v \in \mathcal{T}\). These tests developed by Qu and Yoon
(\citeproc{ref-quyoon2019}{2019}) are also robust-bias-corrected in the
sense of Calonico et al. (\citeproc{ref-calonicoetal2014}{2014}).

To assess the robustness of the results, we will repeat the analyses
with variations in modelling options such as the polynomial degree,
slight changes of the optimal common bandwidth and different optimal
bandwidths at each side of the cutoff. Additionally, we will also run
subgroup analyses to look for heterogeneous effects. This strategy
repeats the analyses in strata of the population formed by covariates
that might be an effect modifier of the treatment.

\section{Results}\label{sec-results}

\subsection{Effects on assistance take-up}\label{sec-rassistance}

The design of AXL allows us to identify the effect of different
treatment endowments on actual treatment uptake. It seems reasonable to
assume that policymakers give more money to the less-employable people
to encourage them to take a more intense treatment. We can measure the
intensity of such treatment through the number of hours actually spent
in job search assistance (\(Y_{HS}\)).

Figure~\ref{fig-intensity} displays visual evidence on the local means
of hours of treatment uptake for each endowment in function of the score
of the participant. Concerning the first treatment (\(D_1\)), we
appreciate a jump at the cutoff for the actual hours of job search
assistance. Regarding the second treatment (\(D_2\)), the number of
hours of job search assistance also presents a small change of level at
the cutoff, although much smaller than the potential difference allowed
by the endowment.

\begin{figure}

\caption{\label{fig-intensity}RD plots for the effect on assistance
take-up. LSM: Local Sample Mean.}

\begin{minipage}{0.50\linewidth}
\pandocbounded{\includegraphics[keepaspectratio]{../intermediate/script04/plots/late_d1_yhs.pdf}}\end{minipage}%
%
\begin{minipage}{0.50\linewidth}
\pandocbounded{\includegraphics[keepaspectratio]{../intermediate/script04/plots/late_d2_yhs.pdf}}\end{minipage}%

\end{figure}%

To complement the plots with quantitative estimates of the effects,
Table~\ref{tbl-assistance} offers formal estimates for each treatment.
All results are consistent with the graphical evidence. Firstly, we find
a mean difference of only an extra hour of job search assistance in the
active treatment both for \(D_1\) and \(D_2\). This difference is
statistically significant at the 1 \% level in both cases. The magnitude
is remarkably low once we consider that the potential difference between
both groups is 6 hours (in \(D_1\)) and 14 hours (in \(D_2\)). As we
show in the supplementary document, the different amount of the vouchers
for job search assistance does not have a statistically significant
effect on the number of training hours either. This is less surprising
due to the equal amount of money reserved for training in every group.
Therefore, considering the small differences of participation in
assistance between control and active treatments, the results below may
be interpreted as the unique effect of different performance pays for
job search intermediation.

\begin{table}

\caption{\label{tbl-assistance}Estimates of effects for each treatment
assignment at the mean of hours of assistance take-up.}

\centering{

\centering
\begin{talltblr}[         %% tabularray outer open
entry=none,label=none,
note{}={Note: * $p$ < 0.1, ** $p$ < 0.05, *** $p$ < 0.01. Each cell shows the point estimate of the $LATE$ and its standard error in parenthesis. $n_h$ is the effective sample size. The polynomial degree is always one.},
]                     %% tabularray outer close
{                     %% tabularray inner open
width={0.5\linewidth},
colspec={X[]X[]},
hline{2}={1-2}{solid, black, 0.05em},
hline{1}={1-2}{solid, black, 0.08em},
hline{8}={1-2}{solid, black, 0.08em},
column{2}={}{halign=c},
column{1}={}{halign=l},
}                     %% tabularray inner close
& $Y_{HS}$ \\
$D_1$ & {1.255***\\(0.210)} \\
Bandwidth $(h)$ & {0.046} \\
$n_h$ & {4774} \\
$D_2$ & {1.098***\\(0.255)} \\
Bandwidth $(h)$ & {0.074} \\
$n_h$ & {8013} \\
\end{talltblr}

}

\end{table}%

\subsection{Effects on employment}\label{sec-remploy}

Figure~\ref{fig-emplmean} presents regression discontinuity plots with
local means of the number of days worked at the first and second
semester for the first (score \(S_1\)) and the second treatment (score
\(S_2\)). Plots for the remaining semesters are available in the
supplementary document. We see that there is no jump at the cutoff in
any case, which suggests a null effect of both the first increase
(\(D_1\)) and the second increase in the treatment dose (\(D_2\)).
Although the visual evidence offers a negative evaluation of the
interventions, more formal methods are required to quantify such
differences.

\begin{figure}

\caption{\label{fig-emplmean}RD plots with local means for employment
duration variables. LSM: local sample mean.}

\begin{minipage}{0.50\linewidth}
\pandocbounded{\includegraphics[keepaspectratio]{../intermediate/script04/plots/late_d1_y1.pdf}}\end{minipage}%
%
\begin{minipage}{0.50\linewidth}
\pandocbounded{\includegraphics[keepaspectratio]{../intermediate/script04/plots/late_d1_y2.pdf}}\end{minipage}%
\newline
\begin{minipage}{0.50\linewidth}
\pandocbounded{\includegraphics[keepaspectratio]{../intermediate/script04/plots/late_d2_y1.pdf}}\end{minipage}%
%
\begin{minipage}{0.50\linewidth}
\pandocbounded{\includegraphics[keepaspectratio]{../intermediate/script04/plots/late_d2_y2.pdf}}\end{minipage}%

\end{figure}%

Table~\ref{tbl-emplmean} offers formal estimates of the \(LATE\) for
both treatments, with results that are consistent with those presented
in previous plots. The first panel shows the mean effect of \(D_1\) at
the cutoff across the four semesters. We do not have evidence of any
statistically significant average effect at any of the semesters
analysed. Setting aside statistical inference, the point estimates are
uniformly small, ranging from an average decrease of one to an average
increase of eight days worked. The second panel presents the mean effect
of \(D_2\) at the cutoff for the same periods. Considering that the
difference between active and control treatments is larger in this case,
the impact would be larger under a linear dose--response function. Our
findings indicate that the effect is only statistically significant at
the second semester (for \(\alpha=0.05\)), with a moderate effect
magnitude of ten days worked. During the rest of the semesters, the
effects are not statistically significant at the conventional
significance level. Therefore, at least for those at the threshold, our
analyses suggest that the effect of both increases in the treatment dose
is virtually null for most semesters\footnote{The results are robust to
  using ``being employed at quarter \(t\)'' as alternative outcome. The
  estimates on the increases in the employment probability are still
  tiny (overall lower than 2 p.p. in absolute terms) and not
  statistically significant in all quarters.}.

\begin{table}

\caption{\label{tbl-emplmean}Estimates of employment duration effects
for each treatment at the mean, semesters 1 to 4.}

\centering{

\centering
\begin{talltblr}[         %% tabularray outer open
entry=none,label=none,
note{}={Note: * $p$ < 0.1, ** $p$ < 0.05, *** $p$ < 0.01. Each cell shows the point estimate of the $LATE$ and its standard error in parenthesis. $n_h$ is the effective sample size. The polynomial degree is always one.},
]                     %% tabularray outer close
{                     %% tabularray inner open
colspec={Q[]Q[]Q[]Q[]Q[]},
hline{2}={1-5}{solid, black, 0.05em},
hline{1}={1-5}{solid, black, 0.08em},
hline{8}={1-5}{solid, black, 0.08em},
column{2-5}={}{halign=c},
column{1}={}{halign=l},
}                     %% tabularray inner close
& $Y_1$ & $Y_2$ & $Y_3$ & $Y_4$ \\
$D_1$ & {-1.890\\(3.698)} & {4.739\\(4.183)} & {7.961\\(5.115)} & {4.428\\(4.743)} \\
Bandwidth $(h)$ & {0.064} & {0.061} & {0.045} & {0.053} \\
$n_h$ & {6634} & {6375} & {4716} & {5495} \\
$D_2$ & {6.980*\\(4.383)} & {10.141**\\(4.867)} & {4.229\\(4.752)} & {7.178\\(5.256)} \\
Bandwidth $(h)$ & {0.037} & {0.044} & {0.049} & {0.044} \\
$n_h$ & {4077} & {4885} & {5309} & {4823} \\
\end{talltblr}

}

\end{table}%

Even if the treatments caused a null effect at the mean of our outcome
variables, it is possible that they shifted the location of certain
quantiles. Figure~\ref{fig-emplmedian} explores the local medians of
both treatments before and after the threshold for semesters 1 and 2.
Figures for the rest of semesters and for the effects at other quantiles
are available in the supplementary document. For the first active
treatment (\(D_1\)), we barely appreciate a discontinuity at the cutoff
for the number of days worked at semester one or two. For the second
active treatment (\(D_2\)), the situation is quite similar if we focus
on the first semester. The figure for the second semester presents a
more pronounced jump at the median of days worked just at the threshold.
This might be suggesting a positive effect of this treatment on the
employment duration.

\begin{figure}

\caption{\label{fig-emplmedian}RD plots with local medians for
employment duration variables. LSMed: local sample median.}

\begin{minipage}{0.50\linewidth}
\pandocbounded{\includegraphics[keepaspectratio]{../intermediate/script04/plots/lqte_d1_y1.pdf}}\end{minipage}%
%
\begin{minipage}{0.50\linewidth}
\pandocbounded{\includegraphics[keepaspectratio]{../intermediate/script04/plots/lqte_d1_y2.pdf}}\end{minipage}%
\newline
\begin{minipage}{0.50\linewidth}
\pandocbounded{\includegraphics[keepaspectratio]{../intermediate/script04/plots/lqte_d2_y1.pdf}}\end{minipage}%
%
\begin{minipage}{0.50\linewidth}
\pandocbounded{\includegraphics[keepaspectratio]{../intermediate/script04/plots/lqte_d2_y2.pdf}}\end{minipage}%

\end{figure}%

Table~\ref{tbl-emplquant} shows estimates of the \(LQTE(\tau)\) on the
employment duration of each semester, that is, the effect on the
different quartiles of the distribution of days worked. As expected, the
results are consistent with the visual evidence presented just before
for the median. Regarding the first treatment (\(D_1\)), the first main
result is the point estimate of zero effect on the first quartile for
all semesters. Secondly, we find an increase of approximately 20 days of
employment for the median of the third semester of our sample. Beyond
the point estimate, in the supplementary document, we show that the
confidence intervals for all parameters of all quartiles include the
zero. Thereby, we cannot reject that this effect on the median is simply
due to chance. Still, we can reject the hypothesis of no effect at any
quartile in the case of this third semester\footnote{This positive
  effect at the median contrasts with a smaller and not statistically
  significant effect at the mean of the same quarter. The point
  estimates for the first and third quartile derived from a
  bias-corrected estimator are slightly negative, but not statistically
  significant. The same occurs for more extreme quantiles.
  Unfortunately, this less-biased estimation involves more imprecision.
  We argue that the negative impacts we found would explain the
  median-mean difference. These additional robustness checks overall
  support the main results presented in this article and are available
  upon request.}. The point estimate decreases in the next semester and
we cannot reject the treatment non-significance hypothesis at this
period. In conclusion, all this evidence makes us argue that there was
not a persistent effect of \(D_1\) on the distribution of days worked.

\begin{table}

\caption{\label{tbl-emplquant}Estimates of employment duration effects
for each treatment at quantiles, semesters 1 to 4.}

\centering{

\centering
\begin{talltblr}[         %% tabularray outer open
entry=none,label=none,
note{}={Note: * $p$ < 0.1, ** $p$ < 0.05, *** $p$ < 0.01. For the first three rows of each panel, each cell reports the point estimate of the $LQTE(\tau)$. Hypotheses tests are only carried out for those $\tau$’s for which the estimate of $LQTE(\tau) \ne 0$.},
]                     %% tabularray outer close
{                     %% tabularray inner open
colspec={Q[]Q[]Q[]Q[]Q[]},
hline{2}={1-5}{solid, black, 0.05em},
hline{1}={1-5}{solid, black, 0.08em},
hline{14}={1-5}{solid, black, 0.08em},
column{2-5}={}{halign=c},
cell{1,3-7,9-13}{1}={}{halign=l},
cell{2,8}{1}={c=5}{halign=l},
}                     %% tabularray inner close
& $Y_1$ & $Y_2$ & $Y_3$ & $Y_4$ \\
$D_1$ & $D_1$ & $D_1$ & $D_1$ & $D_1$ \\
$\tau = 0.25$ & 0 & 0 & 0 & 0 \\
$\tau = 0.5$ & -5.072 & 10.555 & 21.91 & 12.322 \\
$\tau = 0.75$ & 3.863 & 1.015 & 1.014 & 1 \\
Treatment non-significance test &  &  & ** &  \\
Treatment homogeneity test &  &  &  &  \\
$D_2$ & $D_2$ & $D_2$ & $D_2$ & $D_2$ \\
$\tau = 0.25$ & 0 & 0 & 0 & 0 \\
$\tau = 0.5$ & 12.825 & 58.384 & -7.183 & -1.388 \\
$\tau = 0.75$ & -1.75 & 6.001 & 0 & 0 \\
Treatment non-significance test &  & * &  & * \\
Treatment homogeneity test &  &  &  &  \\
\end{talltblr}

}

\end{table}%

Concerning the second treatment (\(D_2\)), we find increases of
approximately two months of employment at the median of the distribution
for the second semester. This point estimate is of remarkable size
considering that the range of the outcome is {[}0, 184{]} days.
Unfortunately, the estimation is not very precise and, as shown in the
supplementary document, its confidence interval includes the zero.
Therefore, we cannot reject that the observed difference of quantiles is
due to the sampling. The median effects contrast with point estimates
around zero both for the first and the third quartiles of the
distribution in almost every semester. Considering the statistical
significance of our results, we cannot reject the hypothesis of no
effect in any quartile for the number of days worked if we establish the
conventional significance level (at \(\alpha=0.05\)). Thereby, in our
sample, the most relevant distributional impact of treatment \(D_2\) is
a 58 days effect at the median during the second semester. Nevertheless,
this heterogeneous effect fades out along time.

Since assigning more money to participants does not overall generate a
positive \(LATE\) on the employment duration, maybe it does on the
quality of the employments. Unfortunately, we can only estimate the
impact on the whole interval of two years after treatment started.
Figure~\ref{fig-emplprop} offers visual evidence on this issue regarding
two types of employment contracts as possible outcomes of the
treatments: an open-ended contract (OEC) and a fixed-term contract
longer than 12 months (FTCm12). Plots for contracts of shorter duration
are available in the supplementary document. The main message from the
plots is that there is no clear jump at the cutoff, so there seems to be
no effect on employment quality. We can appreciate small discontinuities
for the OEC outcome for both treatments, but this conclusion would
require a more fine-grained analysis with numeric estimates.

\begin{figure}

\caption{\label{fig-emplprop}RD plots with local proportions for
employment quality variables.}

\begin{minipage}{0.50\linewidth}
\pandocbounded{\includegraphics[keepaspectratio]{../intermediate/script04/plots/latej_d1_yoec.pdf}}\end{minipage}%
%
\begin{minipage}{0.50\linewidth}
\pandocbounded{\includegraphics[keepaspectratio]{../intermediate/script04/plots/latej_d1_yftcm12.pdf}}\end{minipage}%
\newline
\begin{minipage}{0.50\linewidth}
\pandocbounded{\includegraphics[keepaspectratio]{../intermediate/script04/plots/latej_d2_yoec.pdf}}\end{minipage}%
%
\begin{minipage}{0.50\linewidth}
\pandocbounded{\includegraphics[keepaspectratio]{../intermediate/script04/plots/latej_d2_yftcm12.pdf}}\end{minipage}%

\end{figure}%

Table~\ref{tbl-emplquality} presents the formal estimation of the impact
of AXL on the probability of signing each type of contract as the
longest employment contract signed in two years. The columns collect
both the contracts allowed (open-ended or OEC, long temporary or FTCm12
and intermediate temporary or FTC612) and not allowed (short temporary
or FTC05) to receive the payment. The point estimates are overall not
statistically significant at the conventional level (\(\alpha=0.05\)) in
both treatment comparisons. In the comparison of the high
vs.~intermediate endowments (\(D_2\)), we do find an effect of 4.5 p.p.
on the probability of having a short temporary contract during the
longest employment. At the same time, there is a decrease in the
probability of having an open-ended contract of 3.1 p.p. in this
treatment contrast. This might explain the increases at the median of
the distribution that appeared for \(D_2\), because they were of low
magnitude and only at the second semester. Such increase in short
temporary contracts is interesting, because this is not a type of
contract that guarantees the performance-based pay. One possible
interpretation is that caseworkers thought that this short contract
could facilitate getting a longer one that would get them access to the
placement prize.

\begin{table}

\caption{\label{tbl-emplquality}Estimates of employment quality effects
for each treatment and contract of the longest employment spell.}

\centering{

\centering
\begin{talltblr}[         %% tabularray outer open
entry=none,label=none,
note{}={Note: * $p$ < 0.1, ** $p$ < 0.05, *** $p$ < 0.01. Each cell shows the point estimate of the $LATE$ and 95 \% confidence interval in brackets.},
]                     %% tabularray outer close
{                     %% tabularray inner open
colspec={Q[]Q[]Q[]Q[]Q[]},
hline{2}={1-5}{solid, black, 0.05em},
hline{1}={1-5}{solid, black, 0.08em},
hline{4}={1-5}{solid, black, 0.08em},
column{2-5}={}{halign=c},
column{1}={}{halign=l},
}                     %% tabularray inner close
& $Y_{OEC}$ & $Y_{FTCm12}$ & $Y_{FTC612}$ & $Y_{FTC05}$ \\
$D_1$ & {0.019*\\(-0.001, 0.061)} & {0\\(-0.01, 0.006)} & {0.018\\(-0.009, 0.046)} & {-0.02\\(-0.062, 0.01)} \\
$D_2$ & {-0.031**\\(-0.062, -0.004)} & {0.002\\(-0.004, 0.008)} & {-0.006\\(-0.038, 0.015)} & {0.045***\\(0.018, 0.092)} \\
\end{talltblr}

}

\end{table}%

To examine the process of job search intermediation, we may compare the
results about the effects of the participation in each group on two
slightly different outcomes: the longest contract
(Table~\ref{tbl-emplquality}) and the awarded contract (tbl-emplawards).
While the former variable captures the most stable contract signed
within two years from the start of the treatment, the latter variable
refers to the contract linked to a placement prize (if any).

First, being in the intermediate-dose group (\(D_1\)) increases the
probability of getting any of the awarded contracts. It seems that, in
this case, a larger incentive for the provider triggers a larger effort
to support the job search. Thereby, the treatment appears to have
accelerated access to an employment to a (small) fraction of the
participants. Interestingly, the point estimates among the three types
of contracts are close and do not follow a linear function. If we remind
the effects on the longest employment contract in two years presented in
Table~\ref{tbl-emplquality}, we see that the active treatment has null
effects on the probability of each type of contract. Thus, along time,
either (a) the control group closes the gap with the active treatment
group within two years or (b) the active treatment group is unable to
retain the employment. In the latter case, policymakers could
incorporate a job maintenance intervention involving meetings during the
period of employment.

Second, we see that being in the high-dose group (\(D_2\)) has no effect
on the awarded contract, but it does on the probability of having a
short-term temporary job as the longest employment contract. On one
side, we could interpret that people in this group required a longer
timeline to adapt themselves to the labour market because they enter
better occupations after treatment. This idea does not find support in
our data, because the probability of getting a job with the occupation o
during the longest employment contract barely changes among
groups\footnote{Results are available upon request}. On the other side,
we could interpret that caseworkers push treated jobseekers to accept
less stable contracts with the hope of getting access to more stable
contracts in the future. This would explain the null effects shown in
Table~\ref{tbl-emplawards} on the probability of signing any of the
awarded contracts. However, as we can see in
Table~\ref{tbl-emplquality}, participants in the high-dose group do not
achieve a greater probability of signing a longer temporary contract and
decrease their probability of getting a permanent contract.

\begin{table}

\caption{\label{tbl-emplawards}Estimates of employment quality effects
for each treatment and contract of the awarded employment spell.}

\centering{

\centering
\begin{talltblr}[         %% tabularray outer open
entry=none,label=none,
note{}={Note: * $p$ < 0.1, ** $p$ < 0.05, *** $p$ < 0.01. Each cell shows the point estimate of the $LATE$ and 95 \% confidence interval in brackets.},
]                     %% tabularray outer close
{                     %% tabularray inner open
colspec={Q[]Q[]Q[]Q[]},
hline{2}={1-4}{solid, black, 0.05em},
hline{1}={1-4}{solid, black, 0.08em},
hline{4}={1-4}{solid, black, 0.08em},
column{2-4}={}{halign=c},
column{1}={}{halign=l},
}                     %% tabularray inner close
& $Y_{pOEC}$ & $Y_{pFTCm12}$ & $Y_{pFTC612}$ \\
$D_1$ & {0.022***\\(0.008, 0.05)} & {0.01***\\(0.005, 0.021)} & {0.048***\\(0.019, 0.073)} \\
$D_2$ & {0.008\\(-0.019, 0.024)} & {0.001\\(-0.011, 0.009)} & {0.019\\(-0.007, 0.055)} \\
\end{talltblr}

}

\end{table}%

\subsection{Cream-skimming}\label{sec-rcream}

One of the common rationales to establish different performance pays for
each profile of workers is to discourage cream skimming of users,
i.e.~the selection of easy-to-place workers
(\citeproc{ref-carterwhitworth2015}{Carter and Whitworth 2015}).
Policymakers anticipate that possibility and may include this design
feature to prevent it, but there is no guarantee that the incentives are
high enough to avoid this behaviour. In this section, we evaluate
whether such discrimination exists under the current programme design.
Considering that each provider might have different internal policies
about the prioritization of users and that users self-select in a
specific provider, we stratify the sample by provider for the subsequent
analyses.

Figure~\ref{fig-creaming} shows that there is almost no evidence of
cream skimming, since we cannot reject that being in a less-employable
treatment group changes the participation probability at the
cutoff\footnote{To avoid small sample sizes, we only considered the set
  of providers that accounted for 75\% of the sample (having previously
  sorted the providers from more to less users of AXL).}. Departing from
a theoretical model in which the untreated and the treated potential
outcomes are monotonically related
(\citeproc{ref-heckmanetal2002}{Heckman et al. 2002}), providers could
interpret \(S\) (or \(D\)) as a proxy of the untreated potential outcome
and select individuals based on that score. If cream skimming occurred,
we would appreciate a jump in the regression discontinuity plots or,
equivalently, a non-null \(LATE\) on the participation probability when
the endowment (the profile assigned with the score) is higher.

\begin{figure}

\caption{\label{fig-creaming}Estimates of the \(LATE_j\) on the
probability of participating in each provider. Each dot represents the
point estimate and the bars show the 95 \% confidence interval.}

\begin{minipage}{0.50\linewidth}
\pandocbounded{\includegraphics[keepaspectratio]{../intermediate/script04/plots/creaming_d1.pdf}}\end{minipage}%
%
\begin{minipage}{0.50\linewidth}
\pandocbounded{\includegraphics[keepaspectratio]{../intermediate/script04/plots/creaming_d2.pdf}}\end{minipage}%

\end{figure}%

For both treatments, in general, we do not detect a statistically
significant effect of being assigned to a group with larger endowment on
the probability of participating with each provider. Interestingly, two
of the three coefficients that are statistically significant are
positive, which indicates that a greater endowment fostered the
participation in certain providers. This indicates an ``inverse cream
skimming'' in which providers select hard-to-place workers, something
that could interest some governments in terms of fairness. The unique
exception is provider 5 in \(D_2\), for which we estimate that a larger
treatment endowment decreases the probability of participating in that
provider at the threshold. In any case, this evidence of cream skimming
appears incidental rather than systematic, because that provider only
represents around 1 \% of the sample. Therefore, on a positive note, the
programme appears to have effectively deterred discrimination in the
access to job search services.

\subsection{Plausibility and robustness checks}\label{sec-plaurob}

We have performed two plausibility checks to analyse placebo outcomes
and covariate balance. Then, we have carried out some robustness checks
by varying the bandwidth, the kernel, or the polynomial degree. Here we
summarize the main results of both exercises, while in the supplementary
document the reader can find all the estimates.

The first plausibility check employs placebo outcomes to search evidence
about the presence of an alternative programme using the same scoring
variable to assign treatments. Leveraging the staggered implementation
of the AXL treatments, we can estimate pseudo-treatment effects at the
moment the programme was already implemented but certain individuals had
not yet participated. Differences at these pretreatment periods (for
individuals treated at semester \(t \ge 2\)) would inform on possible
violations of \(\{Y_i(1), Y_i(0)\} \perp\!\!\!\perp (\tilde{S}_i)\) for
individuals treated at semester \(t=1\). If we focus on the effects at
the average, the point estimates are small and not statistically
significant at the usual level. This would support the absence of an
alternative treatment that might bias the estimation of the AXL effect
due to a discontinuity in the untreated potential outcome functions.
Additionally, the documents on the programme design inform that the
score function was published in November 2017 and the mandate to create
such function was passed in July 2017
(\citeproc{ref-veneto2017a}{Government of Veneto 2017a}). This finding
goes in line with the idea of no alternative treatment at
\(\tilde{S}_i=c\) before AXL implementation, since other agents did not
know --at least through open channels-- the score function.

The second plausibility check, the covariate balance analyses, inform us
about the comparability of active and control groups to bring evidence
on the unconfoundedness condition. Attending the magnitude of the
differences , only 1 out of the 52 variables/categories analysed (26 for
each treatment comparison) presents an imbalance higher than the 0.1
recommended threshold derived from Stuart et al.
(\citeproc{ref-stuartetal2013}{2013}) and Austin
(\citeproc{ref-austin2009}{2009}). Following Rosenbaum
(\citeproc{ref-rosenbaum2010}{2010}), we also looked at the
unstandardized mean difference to interpret the size of the problem.
There is an average difference of one year of age between groups for
\(D_1\), which corresponds to 0.13 standard deviations. Considering the
low standardized difference and the practical significance of one year
for the employment probability during most of the working life, we argue
that the confounding effect of the tiny difference of age is negligible
for our conclusions\footnote{As an additional robustness check, we rerun
  the \(LATE\) analyses controlling for the residual imbalance in those
  covariates with differences of at least 5 p.p. (or 0.05 standard
  deviations) and conclusions do not change. Results are available upon
  request.}.

The robustness checks deal with certain changes in the modelling
decisions and their effects on the conclusions. The first set of
analyses preserve the same population of interest. Regarding the
\(LATE\) results, we have varied two modelling options: bandwidth
selector (MSE-optimal but allowing different bandwidths at each side of
the cutoff and a slightly lower common MSE-optimal bandwidth) and
polynomial degree (1 or 2). Then, we have compared the estimates
presented in Section~\ref{sec-remploy} with the alternative results. Our
results on the \(LATE\) are quite robust to the different specifications
both in magnitude and direction. With the alternative bandwidths, the
results are almost identical up to the differences in statistical
significance that might be explained by the distinct power of tests with
different sample sizes. Concerning the \(LQTE(\tau)\) and \(LATE_j\)
results, the robustness analyses are less established because the
estimators have been recently proposed. For the \(LQTE(\tau)\),
following Qu and Yoon (\citeproc{ref-quyoon2019}{2019}), we have
re-estimated the models with a new bandwidth selector by choosing the
second most conservative bandwidth, rather than the most conservative
(i.e., the smallest). For both comparisons, the results for the median
are qualitatively equal with only small differences in quantitative
terms. In all cases the conclusions are preserved. For the \(LATE_j\),
we have applied a different bandwidth for each category of the outcome
by transforming the multinomial model to a set of binomial models. The
qualitative results are preserved and the quantitative estimates are
almost equal, so the conclusions are robust.

\section{Conclusion}\label{sec-conc}

In this article, we have studied how the intensity of job search
services influences its employment effects. To do so, we have leveraged
the discontinuity design of an innovative policy that combines two
design features for the private provision of public services: vouchers
(for job search assistance) and performance-based payments (for job
search intermediation). In this policy, three treatment groups shared
the access to adult training, but had different available doses of other
components: group \(A\) (with 7 hours of assistance and 1,500 € for
intermediation), group \(B\) (with 13 hours of assistance and 2,000 €
for intermediation), and group \(C\) (with 27 hours of assistance and
3,000 € for intermediation).

The results indicate that the local effect on the average of days worked
was null for all semesters in the first comparison (\(B\) vs.~\(A\)) and
for three out of four semesters in the second comparison (\(C\)
vs.~\(B\)). In this last case, we did detect a small but positive effect
during the second semester that later vanishes. We did not find effects
at the first and third quartile of the distribution of days worked, but
some changes appeared at the median. For the intermediate-low
comparison, we detected an approximate increase of around two months at
the median of the third semester. The high-intermediate comparison
presented a rise of twenty days at the second semester after treatment
started. Nonetheless, both effects disappear at the following semesters.
Lastly, it seems that the different treatment intensities generated
certain impacts on the employment quality, although these contracts did
not trigger longer employment episodes. We found a small positive effect
on the probability of having short temporary contracts as the longest
employment for the high-intermediate contrast, whereas the
intermediate-low contrast presented an increase of permanent contracts
during programme implementation. We also found two important results
regarding implementation. First, the mechanism that might explain these
minor employment impacts did not seem to be job search assistance due to
voucher non-redemption, but job search intermediation. Second, we did
not detect evidence of cream-skimming or discrimination of hard-to-place
jobseekers in the access to job search services under the different
performant payments, which suggests that it is a promising design
feature.

Our findings have several policy implications. Regarding job search
assistance, if policymakers wish to ensure that less-employable
individuals receive more treatment, we have seen that these vouchers do
not ensure it. In following versions of AXL, vouchers might be designed
as those usually available in education policies. In these cases, the
individual may choose the provider, but not the number of hours of
education. Regarding job search intermediation, the payment
differentials studied here do not appear sufficient to generate
sustained employment gains. Nonetheless, it is possible that the
differences between the payments associated with each group were not
high or directed enough to cause the desired increase in caseworkers'
effort.

The limitations of our research are mainly two. Firstly, regression
discontinuity only identifies the effects for individuals just at the
threshold of the score. If the impact of the programme was heterogeneous
as a function of the running variable, our results would not be
generalizable to individuals far from the cutoffs. Secondly, the
flexible expenditure of the voucher allowed that different endowments
did not materialize in remarkably different amounts of assistance
received. On one side, this might induce us to think the results
exclusively as the effects of different payment schemes for
intermediation. On the other side, we could have a more complete view
and interpret the findings as intention-to-treat effects. This last
approach understands the lack of voucher redemption as a rejection of
the assumptions behind the programme design formulated by the
policymakers.

Future research might reinforce either the impact or the implementation
analyses of job search services. On the implementation, it would be
necessary to investigate in greater depth the specific actions carried
out as job search assistance. As remarked by Hernán and Robins
(\citeproc{ref-hernanrobins2023}{2023}), it is required to define in a
precise way the interventions we want to analyze. Nowadays, the public
employment service lacks detailed information on the actions carried out
by each specific provider. On the impact, it would be interesting to
consider control groups made up by individuals that received no
treatment. Although identification would require harder assumptions, the
results would complement the view on the worthiness of the programme.
Lastly, aggregate effects seem quite small, so cost-benefit analyses
might help to decide the pertinence of this programme.

In sum, our work does not find support to reserve more money to
treatments of different intensity since it does not seem to generate
lasting effects on the treated people. Future policy designs might
consider devoting such money to treating more jobseekers instead of
increasing the treatment endowment of each participant.

\section*{Statements and
declarations}\label{statements-and-declarations}
\addcontentsline{toc}{section}{Statements and declarations}

\textbf{Funding}. The research has been funded by the Ministry of
Universities of the Government of Spain through the Plan de Formación
del Profesorado Universitario with reference FPU19/05291. We also
recognize AI4S fellowship within the ``Generación D'' initiative by
Red.es, Ministerio para la Transformación Digital y de la Función
Pública, for talent attraction (C005/24-ED CV1), funded by
NextGenerationEU through PRTR.

\textbf{Competing interests}. The author has no competing interests to
declare that are relevant to the content of this article.

\textbf{Data and computer code availability}. The data that support the
findings of this study are available from Veneto Lavoro, but
restrictions apply to the availability of these data, which were used
under license for the current study, and so are not publicly available.
Please, send an e-mail to Veneto Lavoro
(osservatorio.mdl@venetolavoro.it) for further details about how other
researchers can obtain the data and the procedure to request the data.
All the information needed to proceed from the raw data to the results
of the paper (including code) is available in a
\href{https://github.com/ajunquera/moremoney}{GitHub repository}.

\section*{References}\label{references}
\addcontentsline{toc}{section}{References}

\phantomsection\label{refs}
\begin{CSLReferences}{1}{1}
\bibitem[\citeproctext]{ref-aakviketal2005}
Aakvik A, Heckman JJ, Vytlacil EJ (2005) Estimating treatment effects
for discrete outcomes when responses to treatment vary: An application
to {Norwegian} vocational rehabilitation programs. Journal of
Econometrics 125:15--51.
https://doi.org/\href{https://doi.org/10.1016/j.jeconom.2004.04.002}{10.1016/j.jeconom.2004.04.002}

\bibitem[\citeproctext]{ref-abadiecattaneo2018}
Abadie A, Cattaneo MD (2018) Econometric {Methods} for {Program
Evaluation}. Annual Review of Economics 10:465--503.
https://doi.org/\href{https://doi.org/10.1146/annurev-economics-080217-053402}{10.1146/annurev-economics-080217-053402}

\bibitem[\citeproctext]{ref-austin2009}
Austin PC (2009) Using the {Standardized Difference} to {Compare} the
{Prevalence} of a {Binary Variable Between Two Groups} in {Observational
Research}. Communications in Statistics - Simulation and Computation
38:1228--1234.
https://doi.org/\href{https://doi.org/10.1080/03610910902859574}{10.1080/03610910902859574}

\bibitem[\citeproctext]{ref-barnow2009}
Barnow BS (2009) Vouchers in {U}.{S}. Vocational training programs: An
overview of what we have learned. Zeitschrift für ArbeitsmarktForschung
42:71--84.
https://doi.org/\href{https://doi.org/10.1007/s12651-009-0007-9}{10.1007/s12651-009-0007-9}

\bibitem[\citeproctext]{ref-brenninkblonk2012}
Brenninkmeijer V, Blonk RWB (2012) The effectiveness of the {JOBS}
program among the long-term unemployed: A randomized experiment in the
{Netherlands}. Health Promotion International 27:220--229.
https://doi.org/\href{https://doi.org/10.1093/heapro/dar033}{10.1093/heapro/dar033}

\bibitem[\citeproctext]{ref-burgessetal2017}
Burgess S, Propper C, Ratto M, Tominey E (2017) Incentives in the
{Public Sector}: {Evidence} from a {Government Agency}. The Economic
Journal 127:F117--F141.
https://doi.org/\href{https://doi.org/10.1111/ecoj.12422}{10.1111/ecoj.12422}

\bibitem[\citeproctext]{ref-calonicoetal2017}
Calonico S, Cattaneo MD, Farrell MH, Titiunik R (2017) Rdrobust:
{Software} for {Regression-discontinuity Designs}. The Stata Journal:
Promoting communications on statistics and Stata 17:372--404.
https://doi.org/\href{https://doi.org/10.1177/1536867X1701700208}{10.1177/1536867X1701700208}

\bibitem[\citeproctext]{ref-calonicoetal2014}
Calonico S, Cattaneo MD, Titiunik R (2014) Robust {Nonparametric
Confidence Intervals} for {Regression-Discontinuity Designs}: {Robust
Nonparametric Confidence Intervals}. Econometrica 82:2295--2326.
https://doi.org/\href{https://doi.org/10.3982/ECTA11757}{10.3982/ECTA11757}

\bibitem[\citeproctext]{ref-cardetal2018}
Card D, Kluve J, Weber A (2018) What {Works}? {A Meta Analysis} of
{Recent Active Labor Market Program Evaluations}. Journal of the
European Economic Association 16:894--931.
https://doi.org/\href{https://doi.org/10.1093/jeea/jvx028}{10.1093/jeea/jvx028}

\bibitem[\citeproctext]{ref-carterwhitworth2015}
Carter E, Whitworth A (2015) Creaming and {Parking} in {Quasi-Marketised
Welfare-to-Work Schemes}: {Designed Out Of} or {Designed In} to the {UK
Work Programme}? Journal of Social Policy 44:277--296.
https://doi.org/\href{https://doi.org/10.1017/S0047279414000841}{10.1017/S0047279414000841}

\bibitem[\citeproctext]{ref-cattaneoetal2024}
Cattaneo MD, Crump RK, Farrell MH, Feng Y (2024) On {Binscatter}.
American Economic Review 114:1488--1514.
https://doi.org/\href{https://doi.org/10.1257/aer.20221576}{10.1257/aer.20221576}

\bibitem[\citeproctext]{ref-cattaneoetal2019}
Cattaneo MD, Idrobo N, Titiunik R (2019)
\href{https://doi.org/10.1017/9781108684606}{A {Practical Introduction}
to {Regression Discontinuity Designs}: {Foundations}}, 1st edn.
Cambridge University Press

\bibitem[\citeproctext]{ref-cattaneotitiunik2022}
Cattaneo MD, Titiunik R (2022) Regression {Discontinuity Designs}.
Annual Review of Economics 14:821--851.
https://doi.org/\href{https://doi.org/10.1146/annurev-economics-051520-021409}{10.1146/annurev-economics-051520-021409}

\bibitem[\citeproctext]{ref-cheungetal2025}
Cheung M, Egebark J, Forslund A, Laun L, Rödin M, Vikström J (2025) Does
{Job Search Assistance Reduce Unemployment}? {Evidence} on {Displacement
Effects} and {Mechanisms}. Journal of Labor Economics 43:47--81.
https://doi.org/\href{https://doi.org/10.1086/726384}{10.1086/726384}

\bibitem[\citeproctext]{ref-doerretal2017}
Doerr A, Fitzenberger B, Kruppe T, Paul M, Strittmatter A (2017)
Employment and {Earnings Effects} of {Awarding Training Vouchers} in
{Germany}. ILR Review 70:767--812.
https://doi.org/\href{https://doi.org/10.1177/0019793916660091}{10.1177/0019793916660091}

\bibitem[\citeproctext]{ref-dongshen2018}
Dong Y, Shen S (2018) Testing for {Rank Invariance} or {Similarity} in
{Program Evaluation}. The Review of Economics and Statistics 100:78--85.
https://doi.org/\href{https://doi.org/10.1162/REST_a_00686}{10.1162/REST\_a\_00686}

\bibitem[\citeproctext]{ref-dykeetal2006}
Dyke A, Heinrich CJ, Mueser PR, Troske KR, Jeon K-S (2006) The {Effects}
of {Welfare}-to-{Work Program Activities} on {Labor Market Outcomes}.
Journal of Labor Economics 24:567--607.
https://doi.org/\href{https://doi.org/10.1086/504642}{10.1086/504642}

\bibitem[\citeproctext]{ref-eggersetal2024}
Eggers AC, Tuñón G, Dafoe A (2024) Placebo {Tests} for {Causal
Inference}. American Journal of Political Science 68:1106--1121.
https://doi.org/\href{https://doi.org/10.1111/ajps.12818}{10.1111/ajps.12818}

\bibitem[\citeproctext]{ref-engstrometal2012}
Engström P, Hesselius P, Holmlund B (2012) Vacancy {Referrals}, {Job
Search}, and the {Duration} of {Unemployment}: {A Randomized
Experiment}. LABOUR 26:419--435.
https://doi.org/\href{https://doi.org/10.1111/j.1467-9914.2012.00545.x}{10.1111/j.1467-9914.2012.00545.x}

\bibitem[\citeproctext]{ref-escuderoetal2019}
Escudero V, Kluve J, López Mourelo E, Pignatti C (2019) Active {Labour
Market Programmes} in {Latin America} and the {Caribbean}: {Evidence}
from a {Meta-Analysis}. The Journal of Development Studies
55:2644--2661.
https://doi.org/\href{https://doi.org/10.1080/00220388.2018.1546843}{10.1080/00220388.2018.1546843}

\bibitem[\citeproctext]{ref-fougereetal2009}
Fougère D, Pradel J, Roger M (2009) Does the public employment service
affect search effort and outcomes? European Economic Review 53:846--869.
https://doi.org/\href{https://doi.org/10.1016/j.euroecorev.2009.01.006}{10.1016/j.euroecorev.2009.01.006}

\bibitem[\citeproctext]{ref-frolichsperlich2019}
Frölich M, Sperlich S (2019) Impact evaluation: Treatment effects and
causal analysis. Cambridge University Press, Cambridge

\bibitem[\citeproctext]{ref-veneto2017a}
Government of Veneto (2017a) Direttiva per la sperimentazione
dell'{Assegno} per il lavoro per la ricollocazione di lavoratori
disoccupati

\bibitem[\citeproctext]{ref-veneto2017b}
Government of Veneto (2017b) Modello per la profilazione dei destinatari
dell'{Assegno} per il lavoro

\bibitem[\citeproctext]{ref-veneto2018}
Government of Veneto (2018) Deliberazione della {Giunta Regionale}:
{Punto} 73 dell'{ODG} della seduta del 30/04/2018 ({Estratto} del
verbale)

\bibitem[\citeproctext]{ref-heckmanetal2002}
Heckman JJ, Heinrich C, Smith J (2002) The {Performance} of {Performance
Standards}. The Journal of Human Resources 37:778.
https://doi.org/\href{https://doi.org/10.2307/3069617}{10.2307/3069617}

\bibitem[\citeproctext]{ref-heckmansmith2004}
Heckman JJ, Smith JA (2004) The {Determinants} of {Participation} in a
{Social Program}: {Evidence} from a {Prototypical Job Training Program}.
Journal of Labor Economics 22:243--298.
https://doi.org/\href{https://doi.org/10.1086/381250}{10.1086/381250}

\bibitem[\citeproctext]{ref-hernanrobins2023}
Hernán M, Robins JM (2023) Causal inference: {What} if, First edition.
{Taylor and Francis}, Boca Raton

\bibitem[\citeproctext]{ref-hippwarner2008}
Hipp L, Warner ME (2008) Market {Forces} for the {Unemployed}? {Training
Vouchers} in {Germany} and the {USA}. Social Policy \& Administration
42:77--101.
https://doi.org/\href{https://doi.org/10.1111/j.1467-9515.2007.00589.x}{10.1111/j.1467-9515.2007.00589.x}

\bibitem[\citeproctext]{ref-huberetal2018}
Huber M, Lechner M, Strittmatter A (2018) Direct and {Indirect Effects}
of {Training Vouchers} for the {Unemployed}. Journal of the Royal
Statistical Society Series A: Statistics in Society 181:441--463.
https://doi.org/\href{https://doi.org/10.1111/rssa.12279}{10.1111/rssa.12279}

\bibitem[\citeproctext]{ref-imbenskaly2012}
Imbens G, Kalyanaraman K (2012) Optimal {Bandwidth Choice} for the
{Regression Discontinuity Estimator}. The Review of Economic Studies
79:933--959.
https://doi.org/\href{https://doi.org/10.1093/restud/rdr043}{10.1093/restud/rdr043}

\bibitem[\citeproctext]{ref-kingzen2006}
King G, Zeng L (2006) The {Dangers} of {Extreme Counterfactuals}.
Political Analysis 14:131--159.
https://doi.org/\href{https://doi.org/10.1093/pan/mpj004}{10.1093/pan/mpj004}

\bibitem[\citeproctext]{ref-koningheinrich2013}
Koning P, Heinrich CJ (2013) Cream-{Skimming}, {Parking} and {Other
Intended} and {Unintended Effects} of {High-Powered}, {Performance-Based
Contracts}: {Cream-Skimming}, {Parking} and {Other Intended} and
{Unintended Effects}. Journal of Policy Analysis and Management
32:461--483.
https://doi.org/\href{https://doi.org/10.1002/pam.21695}{10.1002/pam.21695}

\bibitem[\citeproctext]{ref-linfanteetal2019}
Linfante G, Radicchia D, Stocco P, Toti E (2019) Rapporto di valutazione
della sperimentazione dell'assegno di ricollocazione. ANPAL

\bibitem[\citeproctext]{ref-liuetal2014}
Liu S, Huang JL, Wang M (2014) Effectiveness of job search
interventions: {A} meta-analytic review. Psychological Bulletin
140:1009--1041.
https://doi.org/\href{https://doi.org/10.1037/a0035923}{10.1037/a0035923}

\bibitem[\citeproctext]{ref-manskipepper2018}
Manski CF, Pepper JV (2018) How {Do Right-to-Carry Laws Affect Crime
Rates}? {Coping} with {Ambiguity Using Bounded-Variation Assumptions}.
Review of Economics and Statistics 100:232--244.
https://doi.org/\href{https://doi.org/10.1162/REST_a_00689}{10.1162/REST\_a\_00689}

\bibitem[\citeproctext]{ref-martinsonetal2019}
Martinson K, Harvill E, Litwok D, Schwartz D, Rosa SMDL, Saunders C,
Bell S (2019) Implementation and {Relative Impacts} of {Two Job Search
Assistance Programs} in {New York City}. {Office of Planning, Research,
and Evaluation, Administration for Children and Families, U.S.
Department of Health and Human Services}, Washington, DC

\bibitem[\citeproctext]{ref-meyer1995}
Meyer BD (1995) \href{https://www.jstor.org/stable/2728911}{Lessons from
the {U}.{S}. Unemployment insurance experiments}. Journal of Economic
Literature 33:91--131

\bibitem[\citeproctext]{ref-munozdebustilloetal2011}
Muñoz De Bustillo R, Fernández-Macías E, Antón J-I, Esteve F (2011)
\href{https://doi.org/10.4337/9781849805919}{Measuring {More} than
{Money}: {The Social Economics} of {Job Quality}}. Edward Elgar
Publishing

\bibitem[\citeproctext]{ref-peietal2022}
Pei Z, Lee DS, Card D, Weber A (2022) Local {Polynomial Order} in
{Regression Discontinuity Designs}. Journal of Business \& Economic
Statistics 40:1259--1267.
https://doi.org/\href{https://doi.org/10.1080/07350015.2021.1920961}{10.1080/07350015.2021.1920961}

\bibitem[\citeproctext]{ref-quyoon2015}
Qu Z, Yoon J (2015) Nonparametric estimation and inference on
conditional quantile processes. Journal of Econometrics 185:1--19.
https://doi.org/\href{https://doi.org/10.1016/j.jeconom.2014.10.008}{10.1016/j.jeconom.2014.10.008}

\bibitem[\citeproctext]{ref-quyoon2019}
Qu Z, Yoon J (2019) Uniform {Inference} on {Quantile Effects} under
{Sharp Regression Discontinuity Designs}. Journal of Business \&
Economic Statistics 37:625--647.
https://doi.org/\href{https://doi.org/10.1080/07350015.2017.1407323}{10.1080/07350015.2017.1407323}

\bibitem[\citeproctext]{ref-rosenbaum2010}
Rosenbaum PR (2010)
\href{https://doi.org/10.1007/978-1-4419-1213-8}{Design of
{Observational Studies}}. Springer New York, New York, NY

\bibitem[\citeproctext]{ref-rubin2007}
Rubin DB (2007) The design {\emph{versus}} the analysis of observational
studies for causal effects: Parallels with the design of randomized
trials. Statistics in Medicine 26:20--36.
https://doi.org/\href{https://doi.org/10.1002/sim.2739}{10.1002/sim.2739}

\bibitem[\citeproctext]{ref-stuartetal2013}
Stuart EA, Lee BK, Leacy FP (2013) Prognostic score--based balance
measures can be a useful diagnostic for propensity score methods in
comparative effectiveness research. Journal of Clinical Epidemiology
66:S84--S90.e1.
https://doi.org/\href{https://doi.org/10.1016/j.jclinepi.2013.01.013}{10.1016/j.jclinepi.2013.01.013}

\bibitem[\citeproctext]{ref-titiunik2021}
Titiunik R (2021)
\href{https://doi.org/10.1017/9781108777919.008}{Natural {Experiments}}.
In: Druckman J, Green DP (eds) Advances in {Experimental Political
Science}, 1st edn. Cambridge University Press, pp 103--129

\bibitem[\citeproctext]{ref-vandenbergetal2019}
Van Den Berg GJ, Hofmann B, Uhlendorff A (2019) Evaluating {Vacancy
Referrals} and the {Roles} of {Sanctions} and {Sickness Absence}*. The
Economic Journal uez032.
https://doi.org/\href{https://doi.org/10.1093/ej/uez032}{10.1093/ej/uez032}

\bibitem[\citeproctext]{ref-voorenetal2019}
Vooren M, Haelermans C, Groot W, Maassen Van Den Brink H (2019) The
{Effectiveness} of {Active Labor Market Policies}: {A Meta}-{Analysis}.
Journal of Economic Surveys 33:125--149.
https://doi.org/\href{https://doi.org/10.1111/joes.12269}{10.1111/joes.12269}

\bibitem[\citeproctext]{ref-xu2017}
Xu K-L (2017) Regression discontinuity with categorical outcomes.
Journal of Econometrics 201:1--18.
https://doi.org/\href{https://doi.org/10.1016/j.jeconom.2017.07.004}{10.1016/j.jeconom.2017.07.004}

\bibitem[\citeproctext]{ref-zanellasalomone2025}
Zanella G, Salomone R (2025) Incentive-based active labor market
programs: {Insights} from policy experimentation in {Italy}. Labour
Economics 93:102687.
https://doi.org/\href{https://doi.org/10.1016/j.labeco.2025.102687}{10.1016/j.labeco.2025.102687}

\end{CSLReferences}




\end{document}
